\documentclass[10pt,letterpaper]{article}
\usepackage[margin=0.90in]{geometry}
\geometry{left=0.90in,bottom=0.90in}
\usepackage{enumitem}
\usepackage{tabularx}
\usepackage{hyperref}
\usepackage{titlesec}
\usepackage{lmodern}
\usepackage[T1]{fontenc}
\hypersetup{colorlinks=true,urlcolor=black,linkcolor=black}
\setlength{\parindent}{0pt}\setlength{\parskip}{0pt}
\titleformat{\section}{\large\bfseries}{}{0pt}{}[\vspace{-4pt}\rule{\textwidth}{0.4pt}]
\titlespacing*{\section}{0pt}{8pt}{5pt}
\setlist[itemize]{leftmargin=1.15em,itemsep=1.5pt,topsep=1pt,parsep=0pt,partopsep=0pt}
\newcommand{\heading}[2]{\vspace{2pt}\begin{tabularx}{\textwidth}{@{}Xr@{}}\textbf{#1} & \textit{#2}\end{tabularx}\vspace{-0pt}}
\newcommand{\headingwtitle}[3]{\vspace{2pt}\begin{tabularx}{\textwidth}{@{}Xr@{}}\textbf{#1} & #3\\[1pt] \textit{#2} & \end{tabularx}\vspace{-0pt}}

\begin{document}
\begin{center}
{\LARGE \textbf{Kebei Jiang}}\\[3pt]
Bellevue, WA \quad $\cdot$ \quad
\href{mailto:kebei.jiang13@gmail.com}{kebei.jiang13@gmail.com} \quad $\cdot$ \quad
(225) 200-4286 \quad $\cdot$ \quad
\href{https://www.linkedin.com/in/kebeijiang}{linkedin.com/in/kebeijiang}
\end{center}
\vspace{-3pt}

\section*{Profile}
Applied scientist with 10+ years of experience in AI/ML research, quantitative experimentation, and end-to-end systems. Recent work spans LLM/VLM development, agent evaluation, synthetic benchmarks, LLM-as-a-judge calibration, RAG, multimodal document understanding, and model fine-tuning. Experienced in developing evaluation methodologies, designing controlled experiments, building production-oriented AI systems, and solving ambiguous problems across research and applied settings.

\section*{Core Expertise}
\textbf{AI/ML:} LLMs $\cdot$ VLMs $\cdot$ Generative AI $\cdot$ Model Fine-Tuning $\cdot$ RAG $\cdot$ Agentic Systems $\cdot$ Deep Learning $\cdot$ NLP

\textbf{Prompt Engineering \& Optimization:} Prompt Design $\cdot$ Specification-Driven Prompting $\cdot$ Algorithmic Prompt Optimization

\textbf{Evaluation \& Data:} LLM-as-a-Judge Calibration $\cdot$ Systematic Agent Evaluation $\cdot$ Synthetic Data Generation $\cdot$ Evaluation Benchmarks $\cdot$ Experiment Design $\cdot$ Statistical Modeling $\cdot$ Data Analysis

\textbf{ML Frameworks \& Serving:} PyTorch $\cdot$ HF Transformers $\cdot$ Distributed Training $\cdot$ Model Deployment \& Serving $\cdot$ vLLM $\cdot$ spaCy

\textbf{Tools \& Infrastructure:} Python $\cdot$ R $\cdot$ SQL $\cdot$ FastAPI $\cdot$ Docker $\cdot$ Git $\cdot$ Azure $\cdot$ GCP

\section*{Experience}
\headingwtitle{Microsoft}{Senior Applied Scientist}{Jan 2020 -- Sep 2026}
\begin{itemize}
\item Developed and deployed LLM and VLM solutions spanning agentic AI, RAG, knowledge mining, information extraction, prompt engineering, and model fine-tuning.
\item Designed evaluation methodologies and experiments for LLM-based agents and AI systems, including synthetic benchmarks, behavioral diagnostics, quantitative metrics, and regression evaluation.
\item Developed production ML/AI systems in Python, with experience in FastAPI, Docker, automated testing, production monitoring, and collaborative code reviews.
\item Built end-to-end ML and information-processing systems from data preparation and feature engineering through model development, evaluation, and deployment.
\item Led technical development across scientists, PMs, engineers, architects, and customer engineers, including technical design, implementation reviews, technical decisions, and customer engagements.
\end{itemize}

\headingwtitle{Viridien}{Applied Scientist}{Jan 2015 -- Jan 2020}
\begin{itemize}
\item Applied statistical and machine-learning methods to process and model large-scale (TB) seismic data.
\item Developed CNN-based solutions for seismic interpretation, including pixel-level identification of subsurface salt structures.
\end{itemize}

\section*{Selected Projects \& Research}

\textbf{Work Projects}

\heading{Specification-to-Calibration: Synthetic Supervision for LLM Judges}{2026}
\begin{itemize}
\item Designed and evaluated a novel algorithmic approach for automated LLM judge calibration by compiling natural-language specifications into coverage-controlled behavioral contrast suites.
\item Utilized feedback-conditioned and GEPA-based optimization to empirically demonstrate model adaptation, improving position-debiased accuracy by 3.58 pp on LLMBar benchmarks.
\end{itemize}

\heading{Period-Closing Skill --- Stateful, Multi-Turn Agent Evaluation}{2026}
\begin{itemize}
\item Developed an agentic skill for automating period-closing workflows involving due dates, calendars, worker assignments, and task dependencies.
\item Built an evaluation framework with 30+ reusable, parameterized scenario templates covering diverse failure modes; collaborated with the SMEs to validate scenario realism and coverage.
\item Created stateful, multi-turn test environments with pre-existing tasks, calendar constraints, and similarly named workers; simulated user responses to assess agent behavior across interaction turns.
\item Evaluated both final task outcomes and agent execution traces across repeated runs; measured task success rate (TSR), tool-call accuracy, latency, token cost, and consistency.
\end{itemize}

\heading{Account Reconciliation Agent --- Synthetic Benchmark \& Adaptive Preference Monitoring}{2026}
\begin{itemize}
\item Developed an accounting agent deployed to real users to resolve ledger--subledger reconciliation exceptions.
\item Analyzed production telemetry to identify shifts in user behavior and built a proof of concept demonstrating the agent's ability to adapt.
\item Built a synthetic benchmark with hundreds of LLM-generated transaction exceptions and ground-truth resolutions and tool-use paths; evaluated task success, tool-call accuracy, latency, and consistency across five runs.
\item Improved benchmark realism by validating transaction-matching distributions with the PM, substantially improving evaluation results and establishing fixed offline regression tests.
\end{itemize}

\heading{VLM Fine-Tuning --- Taxonomy Generalization \& Shortcut-Learning Diagnostic}{2025}
\begin{itemize}
\item Designed a hierarchical taxonomy spanning hundreds of fashion-item classes and aligned it with GPT-4o's understanding of retail products; generated $\sim$1M synthetic training examples 
\item Fine-tuned a Microsoft Phi-4-multimodal foundation model (5.6B) using PyTorch across multiple GPUs with LoRA, SFT, and DPO. Achieved $\sim$90\% accuracy against GPT-4o ground truths
\item Designed a rigorous, controlled diagnostic to analyze shortcut learning and taxonomy memorization, evaluating the limits of domain adaptation versus base model behavior.
\end{itemize}

\heading{Multimodal RAG --- Financial Insights Copilot}{2024}
\begin{itemize}
\item Built an end-to-end RAG system for content-rich financial PDFs, spanning document parsing, multimodal indexing, embeddings, retrieval, query decomposition, and knowledge-graph exploration; released as part of Microsoft Financial Insights Copilot.
\item Identified table-parsing issues in complicated financial documents and explored normalization strategies for downstream processing; systematically evaluated BM25 configurations and reranking models, improving offline retrieval by $\sim$20\%.
\item Explored separate text/table/image indices, ranking fusion, CLIP/SigLIP cross-modal matching, LlamaIndex query decomposition, and knowledge-graph/tree-based approaches including RAPTOR, KPG and GraphRAG.
\item Fine-tuned a Phi-3 Vision model on public chart datasets for visual question answering.
\end{itemize}

\heading{End-to-End Knowledge Mining --- Multimodal Document Processing}{2022--2023}
\begin{itemize}
\item Led development of an end-to-end knowledge-mining solution for large multimodal document collections in the energy domain using Azure Cognitive Search.
\item Designed and optimized indexing flows and custom AI skills spanning chart detection, document classification, entity recognition, and summarization; led a team of scientists and collaborated with PMs, customer engineers, and customers.
\item Identified unnecessary intermediate image storage as a major document-processing bottleneck and redesigned the flow to detect and index image bounding boxes before cropping directly from raw documents, reducing parsing time by $\sim$80\%.
\item Served AI skills for customer integration and supported technical questions and handoff.
\end{itemize}

\heading{Rule-Based NER with spaCy}{2021}
\begin{itemize}
\item Developed a rule-based entity extraction system for geo-location entities in Oil \& Gas documents, including class identification and post-processing to reduce false positives from overlapping and ambiguous names.
\item Achieved over 80\% accuracy and implemented the solution as a custom skill in Azure Cognitive Search.
\end{itemize}

\heading{Automatic Seismic Salt Interpretation}{2018}
\begin{itemize}
\item Used a U-Net architecture to identify salt structures at pixel level in seismic data.
\item Implemented stratified folding, ELU, and Lovasz Loss to improve precision; work was accepted by the 3rd International Conference on Information System \& Data Mining.
\end{itemize}

\section*{Personal Projects}

\heading{OptiBuild --- Domain-Agnostic Optimization Agent | Personal Project}{2026}
\begin{itemize}
\item Engineered a domain-agnostic multi-agent system that translates natural-language constraints into structured Operations Research models, enabling optimized component selection across arbitrary domains (e.g., PC hardware, drone components).
\item Developed an orchestration pipeline utilizing Gemini and CP-SAT with TOPSIS-based multi-objective ranking to evaluate intent fidelity and systematically solve for Pareto-optimal configurations.
\item Designed a staged modelization and repair loop utilizing Pydantic contracts, deterministic validation protocols, and structured execution traces to handle infeasible models or missing data dynamically.
\item Implemented rigorous automated testing and security controls, incorporating a 127-case offline regression suite, PII redaction layers, and strict safety gates for execution-level LLM queries and formulas.
\end{itemize}

\section*{Selected Technical Learning}
\heading{CS336: Language Modeling from Scratch, Stanford}{2025}
\begin{itemize}
\item Implemented core components of large language models from scratch using PyTorch, including custom tokenization pipelines, Transformers architecture and pre-training data filtering diagnostics.
\item Built and verified low-level architectural layers, including implementing FlashAttention-2 mechanics using Triton kernels.
\item Programmed post-training alignment algorithms from scratch, including implementations of GRPO, DPO, and SFT frameworks.
\end{itemize}

\heading{Neural Networks: Zero to Hero, Andrej Karpathy}{2023}
\begin{itemize}
\item Programmed foundational deep learning architectures from scratch using PyTorch, including implementations of micrograd, Bigram, and Multi-Layer Perceptron (MLP) layers.
\item Implemented a Transformers-based English--Spanish translator from scratch and trained it on the Helsinki-NLP/opus-100 dataset.
\end{itemize}

\heading{Reinforcement Learning Specialization, Coursera}{2023}
\begin{itemize}
\item Completed the specialization to deepen understanding of reinforcement learning motivated by the use of RLHF in LLM development.
\end{itemize}

\section*{Education}
\headingwtitle{Louisiana State University, Baton Rouge, LA}{PhD, Physics}{2007 -- 2014}

\section*{Publications}
\begin{itemize}
\item Guang Yang, Kebei Jiang, David P. Lundquist, Prerna Singh. \textit{Specification-to-Calibration: Synthetic Supervision for LLM Judges.} Manuscript in Preparation.
\item Moschoula Pternea, Prerna Singh, Abir Chakraborty, Yagna Oruganti, Mirco Milletari, Sayli Bapat, Kebei Jiang. \textit{The RL/LLM Taxonomy Tree: Reviewing Synergies Between Reinforcement Learning and Large Language Models.} Journal of Artificial Intelligence Research, 2024. DOI: 10.1613/jair.1.15960.
\item Yu Zeng, Kebei Jiang, Jie Chen. \textit{Automatic Seismic Salt Interpretation with Deep Convolutional Neural Networks.} Proceedings of the 2019 3rd International Conference on Information System \& Data Mining, pp. 16--20.
\end{itemize}

\section*{Patent}
\begin{itemize}
\item Yinheng Li, Kebei Jiang. \textit{Concurrent Labeling of Sequences of Words and Individual Words.} U.S. patent application, MS\# 411939-US-NP, 2022.
\end{itemize}

\end{document}
